\section{Hiding model}
\label{sec:privacy}

The public spend view contains the old root, nullifier, output commitment, new
root, asset, fee, deployment domain, and proof data.  The intended private data
are the input-note opening and credential, tree path, output opening, and
internal algebraic trace.  The formal development proves witness independence
for the finite algebraic masking model and records the comparisons needed to
carry that statement to the deployed byte encoding.

\subsection{Algebraic witness independence}

The masking construction has four components: the main trace, copied
consistency columns, boundary columns, and an auxiliary component that
satisfies the final linear condition.  Uniform masks are sampled from finite
spaces.  Surjective linear maps show that the visible masked values have the
same distribution for any two compatible witnesses.

The concrete sampler operates on whole 32-bit words.  It uses first success
under rejection sampling, with at most sixteen words for each field limb.  The
formal view includes this sampler, the degree-four extension-field packing,
the kernel correction, the downstream encoding, and an arbitrary deterministic
serialization of the resulting view.

\begin{theorem}[Fixed-schedule witness independence]
\label{thm:modeled-hiding}
Fix a public statement, transcript schedule, and two private witnesses that
satisfy the same statement.  Under the finite sampling and linear-algebra
hypotheses of the masking construction, the encoded algebraic views generated
from the two witnesses have identical distributions.
\end{theorem}

The equality is exact in the finite model.  Any deterministic encoding
preserves it.

\subsection{Bounded retry}

The prover may test up to seventeen masking schedules and publish the
least-indexed successful one.  Selection is part of the public distribution,
so the model includes the retry controller.  Lean proves that the released
selector and chosen view remain witness-independent when each releasable
successful branch has the same law for both witnesses.

The implementation interface for this theorem identifies the real entropy source,
word rejection sampler, proof-or-abort event, retry cap, selection rule, and
public release with their modeled counterparts.  Stating retry at this level
also exposes timing or abort behavior that a fixed-schedule statement would
miss.

\subsection{Byte inventory and hybrid statement}

The modeled public encoding consists of a 169-byte instruction, a 40-byte
sealed-proof header, and a 75,358-byte proof body, for 75,567 bytes in total.
Lean proves that every offset belongs to one declared field or semantic class
and that splitting and concatenation are inverse operations.  This establishes
a complete byte inventory for the modeled encoding.

To compare deployed views, the formalization uses a two-sided hybrid.  Each
side passes through the Fiat--Shamir compiler, Rust transcript, SHA-256 calls,
serialization, salted Merkle commitments, Poseidon2 calls, entropy expansion,
rejection sampling, retry, and the final algebraic view.  If the error assigned
to these steps on side \(b\) is \(B_b\), then
\begin{equation}
  \Delta(V_0,V_1)\leq B_0+B_1,
  \label{eq:zk-hybrid}
\end{equation}
where \(\Delta\) is statistical distance and
\cref{thm:modeled-hiding} supplies equality at the center of the hybrid.

Equation~\eqref{eq:zk-hybrid} is a conditional statement about the deployed view.  A
computational hash assumption naturally bounds an adversary's distinguishing
advantage, while statistical distance compares full distributions.  A
concrete implementation theorem must use one metric consistently and prove the
Rust, serialization, entropy, retry, and commitment comparisons in that
metric.  The current checked result is exact witness independence for the
finite masking model together with the explicit hybrid interface above.
